\section{Test Case 4: Three-D simulation in presence of a gas counterflow}

The input in the \textit{input-4} sub-directory repeats the PVP parameters
of Test Case 3 (Sec.~\ref{sec:Test-Case-3:}) and enables the aerodynamic
drag (\texttt{system 4}, \texttt{airdrag yes}) with a gas velocity of
$-1000$ cm/s, opposed to the main jet direction; its random term requires
the stochastic Platen integrator (\texttt{integrator 4}).  Since version
2.0 the case reads its stochastic forcing from the pre-generated Gaussian
pool (Sec.~\ref{sec:Noise-pool}), in the CPU build as in the OpenACC one;
before, it drew a Gaussian block at every step.  The noise realization
therefore differs from earlier records, and the regression baselines of the
case were regenerated.  Seed ensembles over the first 5 million steps (means
over 0.03--0.05 s) show that the physics is unchanged: 8 CPU seeds with the
pool and 8 with the former step-by-step noise agree within about one
standard error for every observable (off-axis distance 2.959 cm, 114.6
active beads, cone angle 21.0 degrees, path length 177.5 cm).

The OpenACC build runs the CPU code, byte for byte, until the jet has 100
beads (about 1.12 million steps, 0.011 s), and then the common Platen device
step; the jet reaches the collector at about 0.015 s and then carries about
115 beads.  The device step inserts a bead in the step
that grows the arrays, as the CPU build does.  In a development build of
version 2.0 the first insertion
after the engagement came one step late, the run left the CPU trajectory
there, and the A30 seed ensemble gave an off-axis distance lower by
$4.5\times10^{-4}$ cm than the CPU one; now 24 A30 seeds and 24 CPU seeds
agree within about one standard error for the off-axis distance, the bead
count, the cone angle and the path length.  One NVIDIA A30 runs the 5
million steps in 539 s, against 1454 s on one CPU core.
